\documentclass[a4paper,11pt]{article}
\usepackage[utf8]{inputenc}
\usepackage[T1]{fontenc}
\usepackage{geometry}
\usepackage{listings}
\usepackage{xcolor}
\usepackage{longtable}
\usepackage{amsmath}
\usepackage{hyperref}

\geometry{margin=1in}

\lstset{
backgroundcolor=\color{gray!10},
basicstyle=\ttfamily\small,
breaklines=false,
 columns=fullflexible,
 keepspaces=true,
frame=single
}

\title{Dokumentacija VyOS usmerjevalnika - startup11}
\author{}
\date{\today}

\begin{document}

\maketitle

\section{Naloga in zasnova omrežja}
V tej dokumentaciji opišemo nalogo podjetja in predlagano omrežno zasnovo. Podjetje je start-up brez obstoječe IT infrastrukture; cilj je postaviti ločena omrežna segmenta za uporabnike in za strežnike, zagotoviti varen dostop do interneta in izpostaviti le izbrane javne storitve.

\subsection{Cilji naloge}
\begin{itemize}
  \item Postaviti omrežje z ločenimi segmenti za \textbf{uporabnike} in \textbf{strežnike} (strogo ločeno prometno okolje).
  \item Vsi strežniki so gostovani v DMZ (po zahtevah naloge) — javne in interne storitve bodo fizično v DMZ, vendar z omejitvami dostopa.
  \item Izpostaviti navzven le izbrane storitve: WireGuard, \texttt{wg-portal} (HTTPS) in REST frontend (HTTPS).
  \item Kritične administrativne in imenik storitve (AD, DNS, SNMP) dostopne samo znotraj omrežja in preko VPN (ni direktnega WAN dostopa).
  \item Ohranjati poseben \textbf{ipv6only} segment za eksperimentalne/izobraževalne potrebe z NPTv6 preslikavo za odhodni promet.
\end{itemize}

\subsection{Opredelitev segmentov}
\begin{itemize}
  \item \textbf{DMZ (eth1, 192.168.X.0/24)}: vsi strežniki (REST, wg-portal, WireGuard endpoint, AD/DNS, SNMP, ostali). Fizicni lokaciji strežnikov sta v DMZ, vendar dostopi do nekaterih storitev omejeni z omrežnimi ACL.
  \item \textbf{INTERNAL (eth2, 10.X.0.0/24)}: uporabniške delovne postaje, administrativne postaje in monitoring hosti. Od tu je dovoljen nadzorovan dostop do AD/DNS/SNMP v DMZ.
  \item \textbf{IPV6ONLY (eth3, ULA + NPTv6)}: eksperimentalni IPv6-only segment; omogoča učenje in testiranje IPv6 funkcionalnosti. Izhodni promet je preko NPTv6 preslikave na javni IPv6.
\end{itemize}

\subsection{Izolacija segmentov}
Segmenti se izolirajo z naslednjimi ukrepi:
\begin{itemize}
  \item \textbf{Layer-3 ločitev (VyOS routing + ACL)}: vsak segment ima ločen subnet in VyOS izvaja routing ter aplikacijo politk (default-deny ter strog seznam izjem).
  \item \textbf{VLAN/PortGroup} na hipervizorju: DMZ, INTERNAL in IPV6ONLY so v ločenih PortGroup/VLANih, management port group pa ločen in omejen.
  \item \textbf{Host-firewall} (vsak strežnik in vsaka delovna postaja): default-deny inbound, dovoljenje samo za nujne storitve; admin dostop omejen na internal in VPN izvore.
  \item \textbf{Split DNS / conditional forward}: notranji klienti uporabljajo AD DNS za `startup11.local`, zunanji resolverji pa javne strežnike — prepreči izpostavitev notranjih zapisov.
  \item \textbf{Monitoring in logging}: monitoring host ima izključno dovoljenje za SNMP/agent promet; centralizirano beleženje (syslog) in netflow/sflow za analizo.
\end{itemize}

\subsection{Načrt požarnega zidu (visoka raven)}
Osnovna politika: \emph{default deny} za vse vhodne in posredovane povezave, dovoli le eksplicitno definirane poti.
\begin{itemize}
  \item \textbf{WAN (eth0) -> VyOS}: dovoli \texttt{established, related}; dovoli DNAT za:
    \begin{itemize}
      \item UDP 51820 -> WireGuard endpoint v DMZ
      \item TCP 443 -> reverse-proxy ali neposredno na \texttt{wg-portal}/REST frontend (DMZ)
    \end{itemize}
  \item \textbf{WAN -> DMZ (direktni)}: ne dovoli neposrednega dostopa do AD/DNS/SNMP; dovoli samo zgoraj navedene DNAT pravila.
  \item \textbf{INTERNAL -> DMZ}: dovoli le potrebne storitve:
    \begin{itemize}
      \item DNS (53) do AD/DNS
      \item LDAP/LDAPS (389/636), Kerberos (88), Global Catalog (3268/3269) do AD
      \item HTTP/HTTPS do REST frontend
      \item SSH/RDP samo iz administrativne podmreže ali VPN
    \end{itemize}
  \item \textbf{DMZ -> INTERNAL}: privzeto blokiraj; dovoli le potrebno (npr. monitoring/reporting na notranji monitoring host), definirano z virnim IP-jem.
  \item \textbf{INTERNAL -> INTERNET}: dovoli outbound HTTP/HTTPS/DNS z NAT/masquerade na VyOS.
  \item \textbf{VPN (WireGuard)}: uporabniški VPN ima omejen dostop; administrativni VPN ima dostop do DMZ in internal za upravljanje.
  \item \textbf{Host-firewall pravila}: na vsakem strežniku implementiraj default-deny inbound; dovoli le servise, ki jih strežnik ponuja, z omejitvijo vira na internal/VPN ali specificirane monitoring host-e.
\end{itemize}

V nadaljevanju dokumentacije so podrobnosti konfiguracij in primeri ukazov za VyOS in predloge za host-firewall.

\section{Osnovne informacije}
\begin{itemize}
    \item \textbf{Ime naprave:} startup11-vyos
    \item \textbf{Domena:} startup11.local
    \item \textbf{Uporabnik:} vyos (prijava prek SSH ključa)
\end{itemize}

\section{Omrežni vmesniki}
\begin{longtable}{|l|l|l|l|}
\hline
\textbf{Vmesnik} & \textbf{Opis} & \textbf{IPv4} & \textbf{IPv6} \\
\hline
eth0 & javni internet & 88.200.24.241/25 & 2001:1470:fffd:a8::2/64 \\
\hline
eth1 & DMZ & 192.168.11.1/24 & 2001:1470:fffd:a9::1/64 \\
\hline
eth2 & interna LAN & 10.11.0.1/24 & 2001:1470:fffd:aa::1/64 \\
\hline
eth3 & IPv6-only & - & fd11:11:11::1/64 \\
\hline
\end{longtable}

\section{Privzete usmeritve (statične usmeritve)}
\begin{longtable}{|l|l|}
\hline
\textbf{Destinacija} & \textbf{Naslednji hop} \\
\hline
0.0.0.0/0 & 88.200.24.129 \\
\hline
::/0 & 2001:1470:fffd:a8::1 \\
\hline
\end{longtable}

\section{NAT (IPv4)}
\begin{longtable}{|l|l|l|}
\hline
\textbf{Vrsta} & \textbf{Viri / kriterij} & \textbf{Prevod / opomba} \\
\hline
Izvorni NAT (SNAT) & 10.11.0.0/24, izhodni vmesnik eth0 & maskiranje (masquerade) \\
\hline
Izvorni NAT (SNAT) & 192.168.11.0/24, izhodni vmesnik eth0 & maskiranje (masquerade) \\
\hline
\end{longtable}

\section{DMZ strežniki (vsi v omrežju 192.168.11.0/24)}

\subsection{Pregled strežnikov}
\begin{longtable}{|l|l|l|l|}
\hline
\textbf{Ime strežnika} & \textbf{IP naslov} & \textbf{MAC naslov} & \textbf{Namen} \\
\hline
raft1 & 192.168.11.101 & 00:0c:29:15:42:6f & Raft (cluster node) \\
\hline
raft2 & 192.168.11.102 & 00:0c:29:2b:64:27 & Raft (cluster node) \\
\hline
raft3 & 192.168.11.103 & 00:0c:29:a2:5c:63 & Raft (cluster node) \\
\hline
rest & 192.168.11.104 & 00:0c:29:17:1a:72 & REST API / ostali servisi \\
\hline
dns-srv & 192.168.11.105 & 00:0c:29:4a:b1:99 & DNS strežnik (lokalni) \\
\hline
wireguard & 192.168.11.106 & 00:0c:29:c4:d1:52 & WireGuard VPN končna točka \\
\hline
snmp & 192.168.11.107 & 00:0c:29:69:75:09 & SNMP / monitoring cilj \\
\hline
new\_wg & 192.168.11.108 & 00:0c:29:f9:2a:a8 & New WireGuard-related host \\
\hline
AD (dmz windows 1) & 192.168.11.201 & 00:0c:29:07:cf:1c & Active Directory (Domain Controller) \\
\hline

\end{longtable}

\section{NAT in preusmeritve vrat}
\subsection{IPv4 DNAT (port forwarding)}
\begin{itemize}
    \item \textbf{Pravila:} zunanja vmesnik \texttt{eth0}, UDP port \texttt{51820} preusmerjen na \texttt{192.168.11.106} (WireGuard).
\end{itemize}
\begin{longtable}{|l|l|l|}
\hline
\textbf{Tip} & \textbf{Kriterij} & \textbf{Cilj / opomba} \\
\hline
DNAT (preusmeritev vrat) & vhod: eth0, proto UDP, dst port 51820 & 192.168.11.106 (WireGuard) \\
\hline
SNAT / maskiranje & 10.11.0.0/24 -> izhod eth0 & maskiranje (masquerade) \\
\hline
SNAT / maskiranje & 192.168.11.0/24 -> izhod eth0 & maskiranje (masquerade) \\
\hline
NAT66 & fd11:11:11::/64 -> izhod eth0 & 2001:1470:fffd:ab::/64 \\
\hline
\end{longtable}

\subsection{IPv4 source NAT (masquerade)}
Masquerade je nastavljen za notranji in DMZ promet, ki gre preko javnega vmesnika \texttt{eth0}.
\begin{lstlisting}
# 10.11.0.0/24 -> eth0 (masquerade)
set nat source rule 100 outbound-interface name 'eth0'
set nat source rule 100 source address '10.11.0.0/24'
set nat source rule 100 translation address 'masquerade'

# 192.168.11.0/24 -> eth0 (masquerade)
set nat source rule 110 outbound-interface name 'eth0'
set nat source rule 110 source address '192.168.11.0/24'
set nat source rule 110 translation address 'masquerade'
\end{lstlisting}

\subsection{NAT66 (NPTv6, IPv6-to-IPv6 Network Prefix Translation - RFC 6296)}
Prevod naslovov: \texttt{fd11:11:11::/64} se prevaja na \texttt{2001:1470:fffd:ab::/64} za promet, ki izhaja prek \texttt{eth0}.

\section{Varnost (firewall)}
\subsection{IPv4 - vhodni promet (input)}
\begin{itemize}
    \item Pravilo za ohranjanje stanj: dovoli \texttt{established} in \texttt{related} povezave.
    \item Pravilo za SSH: na javnem vmesniku \texttt{eth0} je dovoljeno TCP dst port \texttt{22} za novo povezavo.
\end{itemize}
\begin{longtable}{|l|l|p{8cm}|}
\hline
\textbf{IP verzija} & \textbf{Veriga} & \textbf{Ključno pravilo / opis} \\
\hline
IPv4 & INPUT & Dovoli vzpostavljene in sorodne povezave (established, related). \\
\hline
IPv4 & INPUT & Dovoli nove TCP povezave na port 22 na vmesniku eth0 (SSH). \\
\hline
IPv4 & FORWARD & Dovoli nove UDP povezave iz eth0 proti 192.168.11.106:51820 (WireGuard). \\
\hline
IPv6 & INPUT & Dovoli vzpostavljene in sorodne povezave; dovoli nove TCP povezave na port 22 na eth0. \\
\hline
IPv6 & FORWARD & Dovoli nove UDP povezave na port 51820 na vhodnem vmesniku eth0 (WireGuard). \\
\hline
\end{longtable}

\subsection{IPv4 - forward (posredovanje)}
\begin{itemize}
  \item Pravilo za WireGuard: dovoli novo UDP povezavo iz \texttt{eth0} proti \texttt{192.168.11.106:51820} s prehodi na \texttt{eth1}.
  \item Pravila 210-241 eksplicitno dovolijo dostop iz \texttt{eth2} do AD DNS/DC na \texttt{192.168.11.201} (DNS, Kerberos, LDAP, Global Catalog, SMB/RPC, dynamic RPC).
\end{itemize}
% Forward rule summary already in firewall table above.

\subsection{IPv6 - vhodni promet in posredovanje}
IPv6 ima podobna vhodna pravila (dovoli \texttt{established/related} in SSH na \texttt{eth0}), forward pravilo za WireGuard (UDP port 51820 na vhodnem vmesniku \texttt{eth0}) in pravila 210-241 za dostop do AD strežnika na \texttt{2001:1470:fffd:a9::185}.
% IPv6 firewall rules summarized in the firewall table above.

\section{Storitve}
\subsection{DHCPv4}
DHCP server ima dve skupini: \textbf{SERVERS} (192.168.11.0/24) z navedenimi statičnimi mapami in \textbf{USERS} (10.11.0.0/24) za uporabniške naprave z avtomatskim razponom.
\begin{longtable}{|l|l|}
\hline
\textbf{Skupina} & \textbf{Subnet / opomba} \\
\hline
SERVERS & 192.168.11.0/24 (statične mape) \\
\hline
USERS & 10.11.0.0/24 (dinamični range 10.11.0.100-200) \\
\hline
\end{longtable}

\subsection{DHCP - statične mape}
Statične DHCP mape so konfigurirane v DHCP strežniku na VyOS in ujemajo IP naslove s MAC naslovi, kot je prikazano spodaj.
\begin{longtable}{|l|l|l|}
\hline
\textbf{Ime} & \textbf{MAC naslov} & \textbf{IP naslov} \\
\hline
raft1 & 00:0c:29:15:42:6f & 192.168.11.101 \\
\hline
raft2 & 00:0c:29:2b:64:27 & 192.168.11.102 \\
\hline
raft3 & 00:0c:29:a2:5c:63 & 192.168.11.103 \\
\hline
rest & 00:0c:29:17:1a:72 & 192.168.11.104 \\
\hline
dns-srv & 00:0c:29:4a:b1:99 & 192.168.11.105 \\
\hline
wireguard & 00:0c:29:c4:d1:52 & 192.168.11.106 \\
\hline
snmp & 00:0c:29:69:75:09 & 192.168.11.107 \\
\hline
new\_wg & 00:0c:29:f9:2a:a8 & 192.168.11.108 \\
\hline
AD (dmz windows 1) & 00:0c:29:07:cf:1c & 192.168.11.201 \\
\hline
\end{longtable}

\subsection{DHCPv6}
DHCPv6 podeljuje naslove v subnetu \texttt{2001:1470:fffd:a9::/64} z range \texttt{::100} - \texttt{::1ff}. Ime strežnika za DHCPv6 je \texttt{2001:1470:fffd:a9::1}.

\paragraph{Statične DHCPv6 mape}
Spodaj so navedene statične DHCPv6 mape z DUID/identifikatorji, ki se uporabljajo v konfiguraciji VyOS:
\begin{longtable}{|l|l|l|}
\hline
\textbf{Host} & \textbf{DHCPv6 naslov} & \textbf{DUID / identifier} \\
\hline
dmz-windows-AD & 2001:1470:fffd:a9::185 & 00:01:00:01:31:86:a4:b0:00:0c:29:07:cf:26 \\
\hline
raft1 & 2001:1470:fffd:a9::124 & 00:02:00:00:ab:11:2b:e2:d4:90:70:f3:b3:37 \\
\hline
raft2 & 2001:1470:fffd:a9::114 & 00:02:00:00:ab:11:e1:94:d4:8a:18:52:ad:98 \\
\hline
raft3 & 2001:1470:fffd:a9::16b & 00:02:00:00:ab:11:92:37:c8:9a:00:b8:67:77 \\
\hline
rest & 2001:1470:fffd:a9::115 & 00:02:00:00:ab:11:5a:9d:49:65:ce:4a:b2:48 \\
\hline
snmp & 2001:1470:fffd:a9::17a & 00:02:00:00:ab:11:b8:b7:97:e3:eb:3c:a5:52 \\
\hline
wg & 2001:1470:fffd:a9::1c8 & 00:02:00:00:ab:11:80:f9:48:e8:e6:e2:12:59 \\
\hline
new\_wg & 2001:1470:fffd:a9::18d & 00:02:00:00:ab:11:dd:e2:8d:62:21:b3:d9:00 \\
\hline
\end{longtable}

\subsection{DNS (posredovanje in split DNS)}
DNS posredovanje na VyOS sprejema poizvedbe iz notranjih omrežij in uporablja conditional forward za AD domeno.

\textbf{Ključna split-DNS logika:}
\begin{itemize}
  \item poizvedbe za \texttt{startup11.local} se na VyOS preusmerijo na AD DNS \texttt{192.168.11.201};
  \item vse ostale poizvedbe gredo na javne upstream DNS strežnike (ARNES + Cloudflare).
\end{itemize}

\textbf{AD DNS server (authoritative):} \texttt{192.168.11.201} / \texttt{2001:1470:fffd:a9::185}

\subsubsection{Trenutno stanje AD DNS}
Trenutna AD zona \texttt{startup11.local} vsebuje klju\v{c}ne notranje zapise za domeno, DC in DMZ stre\v{z}nike:

\begin{itemize}
  \item korenski zapis \texttt{@} vsebuje A zapis \texttt{192.168.11.201} in AAAA zapis \texttt{2001:1470:fffd:a9::185};
  \item NS in SOA na vrhu cone kažeta na \texttt{win-sgi52j8519e.startup11.local.};
  \item standardni AD SRV zapisi so prisotni za \verb|_gc._tcp|, \verb|_kerberos._tcp|, \verb|_kerberos._udp|, \verb|_ldap._tcp| in \verb|_kpasswd._tcp| ter ustrezne \verb|Default-First-Site-Name| in \verb|ForestDnsZones| / \verb|DomainDnsZones| podcone;
  \item glavni DC gostitelj \texttt{win-sgi52j8519e} ima A zapis \texttt{192.168.11.201} in AAAA zapis \newline
\texttt{2001:1470:fffd:a9::185};
  \item dodatna notranja zapisa sta \texttt{rest.startup11.local} \textrightarrow{} \texttt{192.168.11.104} in \texttt{snmp.startup11.local} \textrightarrow{} \texttt{192.168.11.107}.
\end{itemize}

\textbf{Forwarderji na AD DNS:}
\begin{itemize}
  \item IPv4: \texttt{193.2.1.66}, \texttt{1.1.1.1}
  \item IPv6: \texttt{2001:1470:8000::66}, \texttt{2606:4700:4700::1111}
\end{itemize}

\textbf{Dovoljena omrežja za VyOS DNS forwarding:} 10.11.0.0/24, 192.168.11.0/24, 2001:1470:fffd:aa::/64, 2001:1470:fffd:a9::/64, fd11:11:11::/64.

\textbf{Upstream DNS strežniki:}
\begin{longtable}{|l|l|l|}
\hline
\textbf{Ponudnik} & \textbf{IPv4} & \textbf{IPv6} \\
\hline
\endhead
\hline
\endlastfoot
ARNES & 193.2.1.66 & 2001:1470:8000::66 \\
\hline
Cloudflare & 1.1.1.1 & 2606:4700:4700::1111 \\
\hline
\end{longtable}

\subsubsection{Split DNS (dnsmasq)}
Strežnik \texttt{dns-srv} (192.168.11.105) z \textbf{dnsmasq} je opcijski notranji resolver. Trenutna pravilna konfiguracija uporablja \texttt{server=/.../} in ne \texttt{address=/.../} za AD domeno, da se ohranijo SRV zapisi za AD prijavo.

\textbf{Trenutna konfiguracija:}
\begin{itemize}
    \item dnsmasq posluša na: \texttt{192.168.11.105:53} (DMZ naslov)
    \item \texttt{startup11.local} se posreduje na AD DNS: \texttt{192.168.11.201} in \texttt{2001:1470:fffd:a9::185}
    \item Konfiguracija: \verb|/etc/dnsmasq.d/startup11.conf|
\end{itemize}

\textbf{Pomembno za AD prijavo:} notranji odjemalci morajo za \texttt{startup11.local} dobiti SRV zapise iz AD DNS, ne samo A/AAAA zapisov. Zato je v VyOS nastavljen:
\begin{lstlisting}
set service dns forwarding domain startup11.local name-server 192.168.11.201
set service dns forwarding domain startup11.local recursion-desired
\end{lstlisting}

\textbf{Sprememba konfiguracije:}
\begin{lstlisting}
# Edit configuration file
sudo nano /etc/dnsmasq.d/startup11.conf

# Test configuration
sudo dnsmasq --test

# Restart service
sudo systemctl restart dnsmasq

# Check service status
sudo systemctl status dnsmasq

# Check listening on port 53
ss -ltnup '( sport = :53 )'

# Test DNS response
dig +short _ldap._tcp.dc._msdcs.startup11.local @192.168.11.105 SRV
dig +short example.com @192.168.11.105
\end{lstlisting}

\subsubsection{dnsmasq - aktivna konfiguracija in ponovna postavitev}
Spodaj je preverjena konfiguracija, ki trenutno teče na \texttt{192.168.11.105:53}.

\begin{lstlisting}
# /etc/dnsmasq.d/startup11.conf
# Bind only to the DMZ address to avoid conflicts with systemd-resolved on localhost
listen-address=192.168.11.105
bind-interfaces

# Forward AD domain to AD DNS (preserves SRV records)
server=/startup11.local/192.168.11.201
server=/startup11.local/2001:1470:fffd:a9::185

# Upstream resolvers for non-local domains
server=1.1.1.1
server=193.2.1.66
server=2001:1470:8000::66
server=2606:4700:4700::1111
\end{lstlisting}

Koraki za ponovno postavitev:
\begin{enumerate}
  \item Namesti pakete.
  \begin{lstlisting}
sudo apt update
sudo apt install -y dnsmasq dnsutils
  \end{lstlisting}
  \item Ustvari konfiguracijsko datoteko.
  \begin{lstlisting}
sudo tee /etc/dnsmasq.d/startup11.conf > /dev/null << 'EOF'
# (vsebina kot zgoraj)
EOF
  \end{lstlisting}
  \item Preveri sintakso in ponovno zaženi storitev.
  \begin{lstlisting}
sudo dnsmasq --test
sudo systemctl restart dnsmasq
sudo systemctl enable dnsmasq
  \end{lstlisting}
  \item Preveri delovanje.
  \begin{lstlisting}
ss -ltnup '( sport = :53 )'
dig +short _ldap._tcp.dc._msdcs.startup11.local @192.168.11.105 SRV
dig +short example.com @192.168.11.105
  \end{lstlisting}
\end{enumerate}

\subsubsection{mDNS in systemd-resolved: rešitev za .local}

Na nekaterih Ubuntu gostiteljih storitev \texttt{systemd-resolved} obravnava domeno \texttt{.local} kot mDNS (Avahi). Posledično lahko DNS poizvedbe za \texttt{startup11.local} prejmejo odgovor \texttt{REFUSED} ali pa jih sistem poskuša obdelati z mDNS namesto prek omrežnega DNS resolverja.

Na gostitelju, kjer teče \texttt{wg-portal}, smo uporabili preprost popravek: za vmesnik smo nastavili routing domain in onemogočili mDNS, tako da se poizvedbe za \texttt{startup11.local} pošljejo na DHCP-provided resolver (192.168.11.1) in dosežejo AD DNS.

\begin{lstlisting}
# Nastavi routing domain (tilde pomeni routing domain)
sudo resolvectl domain ens160 ~startup11.local

# Onemogoci mdns na vmesniku
sudo resolvectl mdns ens160 no

# Pocisti predpomnilnik, ce je potrebno
sudo resolvectl flush-caches

# Preizkus
resolvectl query startup11.local
dig +short _ldap._tcp.dc._msdcs.startup11.local @192.168.11.1 SRV
\end{lstlisting}

Ta rešitev ohranja mDNS obnašanje za druge gostitelje, hkrati pa zagotavlja, da poizvedbe za \texttt{startup11.local} dosežejo avtoritativni AD DNS.

\subsection{NTP}
????

\section{Monitoring in SNMP}
Za spremljanje metrik v DMZ uporabljamo preprost monitoring stack: Prometheus (scrape in shranjevanje), \texttt{snmp\_exporter} (pretvornik SNMP→Prometheus) in Grafana (vizualizacija). \texttt{snmp\_exporter} pobira metrike z naprave VyOS prek SNMP v2c in jih izpostavi na svojem HTTP vmesniku, ki ga nato pobira Prometheus.

\subsection{Komponente in porti}
\begin{itemize}
  \item \textbf{snmp\_exporter}: 192.168.11.107:9116 (TCP) - endpoint exporterja, npr. \newline
\texttt{http://192.168.11.107:9116/snmp}
  \item \textbf{Prometheus}: 192.168.11.107:9090 (TCP) - scrape cilj in uporabniški vmesnik
  \item \textbf{Grafana}: 192.168.11.107:3000 (TCP) - nadzorne plošče in vizualizacija
  \item \textbf{VyOS SNMP}: 192.168.11.1:161 (UDP) - SNMP agent na usmerjevalniku (v2c community)
\end{itemize}

\subsection{VyOS SNMP (primer)}
Omogočite SNMP na VyOS in ga vežite na DMZ naslov:
\begin{lstlisting}
set service snmp community startup11 authorization 'ro'
set service snmp contact 'ops@startup11.local'
set service snmp listen-address 192.168.11.1
commit
save
\end{lstlisting}

\subsection{Docker Compose primer (snmp\_exporter + Prometheus + Grafana)}
Namestite to na monitoring gostitelja (192.168.11.107) in prilagodite poti/volumne po potrebi.
\begin{lstlisting}
version: '3'
services:
  snmp\_exporter:
    image: prom/snmp-exporter:latest
    ports:
      - "9116:9116"
    volumes:
      - ./snmp.yml:/etc/snmp\_exporter/snmp.yml

  prometheus:
    image: prom/prometheus:latest
    ports:
      - "9090:9090"
    volumes:
      - ./prometheus.yml:/etc/prometheus/prometheus.yml
    depends_on:
      - snmp\_exporter

  grafana:
    image: grafana/grafana:latest
    ports:
      - "3000:3000"
    depends_on:
      - prometheus
\end{lstlisting}

\subsection{Konfiguracija Prometheusa (primer scrape job)}
Dodajte `scrape` nalogo, da Prometheus pobira metrike od exporterja (če ne uporabljate Docker Compose DNS, zamenjajte `snmp\_exporter:9116` z `192.168.11.107:9116`):
\begin{lstlisting}
scrape_configs:
  - job_name: 'snmp'
    static_configs:
      - targets: ['192.168.11.1']
    metrics_path: /snmp
    params:
      module: [vyos]
      auth: [vyos_auth]
    relabel_configs:
      - source_labels: [__address__]
        target_label: __param_target
      - source_labels: [__param_target']
        target_label: instance
      - target_label: __address__
        replacement: 192.168.11.107:9116
\end{lstlisting}

\subsection{Generator in snmp.yml}
Uporabite uradni generator za `snmp\_exporter`, da ustvarite prilagojen `snmp.yml` za VyOS. Postopek v grobem:
\begin{itemize}
  \item Naredite `generator.yml` z `auth` sekcijo (community `startup11`) in modulom `vyos`, ki naredi `walk` za `1.3.6.1.2.1.1.3` (sysUpTime), `1.3.6.1.2.1.2` (ifTable) in `1.3.6.1.2.1.31.1.1` (ifXTable).
  \item Zaženite generator (kot container ali binarko) in dobljeni `snmp.yml` kopirajte na exporter gostitelja v `/etc/snmp\_exporter/snmp.yml`.
\end{itemize}

\subsection{Testiranje}
\begin{lstlisting}
curl "http://192.168.11.107:9116/snmp?module=vyos&target=192.168.11.1"
snmpwalk -v2c -c startup11 192.168.11.1 1.3.6.1.2.1.1.3.0
\end{lstlisting}

\subsection{Router Advertisements (RA)}
RA so omogočeni na:
\begin{itemize}
    \item \texttt{eth1}: prefix \texttt{2001:1470:fffd:a9::/64} (managed flag, no-autonomous)
    \item \texttt{eth2}: prefix \texttt{2001:1470:fffd:aa::/64}
    \item \texttt{eth3}: prefix \texttt{fd11:11:11::/64}, name-server \texttt{fd11:11:11::1}
\end{itemize}

\subsection{SSH}
SSH je konfiguriran za prijavo prek ključev (avtentikacija z geslom je onemogočena) na privzetem portu 22. Če želite začasno omogočiti prijavo z geslom, uredite datoteko \newline  \verb|/etc/ssh/sshd_config.d/50-cloud-init.conf| in nastavite \texttt{PasswordAuthentication yes}, nato ponovno zaženite \texttt{sshd}.

\section{Sistem}
\subsection{Uporabniki in ključi}
Glavni lokalni uporabnik na napravi je \texttt{vyos}. V konfiguraciji so vpisani javni ključi za uporabnika \texttt{vyos}:
\begin{itemize}
    \item \textbf{pavle}: \texttt{AAAAC3NzaC1lZDI1NTE5AAAAIClsTiUna0lG4FgaZOZ8cpxWvlWM7h9B2dbL53+QXr3h}
    \item \textbf{zanzan}: \texttt{AAAAC3NzaC1lZDI1NTE5AAAAIIGRVCofZkoaEaAcW0LquGsYgpRyHZ4Dg0fqVlssZUIL}
\end{itemize}

Za začasno ponovno omogočanje SSH gesel: uredite \verb|/etc/ssh/sshd_config.d/50-cloud-init.conf| in nastavite \texttt{PasswordAuthentication yes}.

\subsection{Uporabniška imena}
Uporabniki na dmz virtualkag so vsi poimenovani \texttt{zanzan0X}, kjer je X številka virtualke.
Na klientih TODO.

\subsection{Splošno geslo}
\begin{center}
\boxed{\texttt{123NajboljseGeslo456!}}
\end{center}

\section{Upravljanje s konfiguracijami}

\subsection{Nalaganje celotne konfiguracije}
Če želite naložiti celotno konfiguracijo iz datoteke (npr. \verb|v4.conf|), uporabite:
\begin{lstlisting}
configure
load v4.conf
commit
save
\end{lstlisting}

\subsection{Apliciranje seznama ukazov}
Če želite aplicirati seznam ukazov iz datoteke (npr. \verb|v4.commands|), uporabite:
\begin{lstlisting}
configure
source v4.commands
commit
save
\end{lstlisting}

\subsection{Opozorilo pri uporabi}
Pri apliciranju ukazov je potrebna pazljivost — nekateri ukazi ne delujejo kot pričakovano:
\begin{itemize}
    \item Ukaz \verb|set| pri vmesnikih (\verb|interface|) ne nastavlja vrednosti, ampak jo dodaja (deluje kot \verb|add|). To lahko vodi do podvajanja nastavitev.
  \\item Originalne konfiguracije \\verb|v1.commands| in \\verb|v3.commands| so zgodovinski osnutki; za trenutno stanje uporabljaj \\verb|v4.conf| in \\verb|v4.commands|.
  \item Najnovejšo in pravilno konfiguracijo (\verb|v4.commands|) generirajte avtomatsko z ukazom (zunaj \verb|configure| moda):
\end{itemize}
\begin{lstlisting}
show configuration commands > v4.commands
\end{lstlisting}

Uporabite to avtomatsko generirano datoteko za zanesljivo apliciranje vseh trenutnih nastavitev.

\section{Active Directory (Windows Server)}
\subsection{Osnovne informacije}
Active Directory domena je nastavljena na strežniku \textbf{sk11-dmz-windows}. Strežnik deluje kot primarni Domain Controller (PDC) za domeno \texttt{startup11.local}.

\begin{itemize}
    \item \textbf{Ime strežnika:} sk11-dmz-windows
    \item \textbf{Domena:} startup11.local
    \item \textbf{NetBIOS ime:} startup11
    \item \textbf{OS:} Windows Server (Core)
    \item \textbf{Vloga:} Domain Controller
\end{itemize}

\subsection{Ustvarjanje uporabnikov v Active Directory}
Uporabnike v domeni ustvarjamo s PowerShell ukazi neposredno na Domain Controllerju.

\textbf{Primer: Ustvarjanje uporabnika}
\begin{lstlisting}[language=bash, breaklines=true]
$password = Read-Host "Enter password for User" -AsSecureString
New-ADUser -Name "Ime Priimek" `
           -GivenName Ime `
           -Surname Priimek `
           -SamAccountName imepriimek `
           -UserPrincipalName ime.priimek@startup11.local `
           -Enabled $true `
           -AccountPassword $password `
           -PassThru
\end{lstlisting}

\subsection{Upravljanje uporabnikov}
\textbf{Prikaz vseh uporabnikov v domeni:}
\begin{lstlisting}[language=bash, breaklines=true]
Get-ADUser -Filter * | Select-Object Name, SamAccountName, Enabled
\end{lstlisting}

\textbf{Dodajanje grupe in dodajanje uporabnika grupi:}
\begin{lstlisting}
New-ADGroup -Name "Group name" -GroupScope Global -GroupCategory Security
Add-ADGroupMember -Identity "Group name" -Members "account name"
\end{lstlisting}

\textbf{Onemogočanje uporabniškega računa:}
\begin{lstlisting}[language=bash, breaklines=true]
Disable-ADAccount -Identity "ime"
\end{lstlisting}

\textbf{Brisanje uporabniškega računa:}
\begin{lstlisting}[language=bash, breaklines=true]
Remove-ADUser -Identity "ime" -Confirm:$false
\end{lstlisting}

\subsection{Prijava v domeno}
Računalnike (Windows Pro/Enterprise ali Linux s konfiguriranim SSSD) lahko priklopimo v domeno \texttt{startup11.local}. Po uspešnem priklopu se uporabniki prijavljajo z:

\begin{itemize}
    \item \textbf{Uporabniško ime:} \texttt{startup11\textbackslash ime} (npr. \texttt{startup11\textbackslash jdoe})
    \item \textbf{Geslo:} (geslo, nastavljeno ob ustvarjanju uporabnika)
\end{itemize}

\subsection{Dodatne informacije}
\begin{itemize}
    \item Administrator domene: \texttt{startup11\textbackslash Administrator} z enakim geslom kot lokalni Administrator na strežniku
\end{itemize}

\subsection{Upravljanje AD DNS streznika}
AD DNS je avtoritativen za \texttt{startup11.local} in vraca AD SRV zapise (npr. \verb|_ldap._tcp.dc._msdcs.startup11.local|).

\textbf{Preverjanje forwarderjev:}
\begin{lstlisting}[language=bash, breaklines=true]
Get-DnsServerForwarder
\end{lstlisting}

\textbf{Preverjanje, da je DNS na AD strezniku lokalen:}
\begin{lstlisting}[language=bash, breaklines=true]
Get-DnsClientServerAddress
Get-DnsServerZone
\end{lstlisting}

\textbf{Izpis DNS zapisov v AD zoni:}
\begin{lstlisting}[language=bash, breaklines=true]
Get-DnsServerResourceRecord -ZoneName "startup11.local" |
  Select-Object HostName, RecordType, RecordData
\end{lstlisting}

\textbf{Preverjanje SRV zapisov na AD DNS:}
\begin{lstlisting}[language=bash, breaklines=true]
Resolve-DnsName _ldap._tcp.dc._msdcs.startup11.local -Type SRV -Server 127.0.0.1
nslookup -type=SRV _ldap._tcp.dc._msdcs.startup11.local 192.168.11.201
\end{lstlisting}

\textbf{Dodajanje notranjih zapisov za DMZ streznike (primer):}
\begin{lstlisting}[language=bash, breaklines=true]
Add-DnsServerResourceRecordA -ZoneName "startup11.local" -Name "rest" -IPv4Address "192.168.11.104"
Add-DnsServerResourceRecordA -ZoneName "startup11.local" -Name "snmp" -IPv4Address "192.168.11.107"
\end{lstlisting}

\textbf{Test novih zapisov po dodajanju:}
\begin{lstlisting}[language=bash, breaklines=true]
Resolve-DnsName rest.startup11.local -Type A -Server 127.0.0.1
Resolve-DnsName snmp.startup11.local -Type A -Server 127.0.0.1
\end{lstlisting}

\textbf{Opomba o split DNS:} VyOS izvaja DNS posredovanje in conditional forward na AD DNS za \texttt{startup11.local}, AD DNS pa ostane avtoritativni vir notranjih zapisov.


\section{WireGuard}

\subsection{Splošne informacije}
WireGuard VPN je nastavljeno na napravi na IP naslovu \texttt{192.168.11.106} v DMZ omrežju. Konfiguracija je bila avtomatsko generirana s \textbf{pivpn} skripto, ki poenostavi upravljanje in vzpostavitev WireGuard VPN strežnika.

Trenutna postavitev ostane nespremenjena v prvi fazi. Podroben fazni načrt, ki ohrani obstoječi PiVPN/NAT in nato v drugi fazi preseli dual WireGuard stack na novo VM z \texttt{wg-portal}, je opisan v \texttt{phased-firewall-plan.md}. Ta novi načrt nadomešča prejšnje pomožne predloge za \texttt{wg-client}.

\subsection{Dostop in port}
\begin{itemize}
    \item \textbf{Notranji naslov:} 192.168.11.106
    \item \textbf{Javni port:} UDP 51820 (preusmeran prek NAT/DNAT na eth0)
    \item \textbf{Storitev:} Aktivna in dostopna iz javnega interneta
\end{itemize}

\subsection{Upravljanje VPN uporabnikov}
WireGuard uporabnike upravljate s pivpn ukazom. Osnovni koraki:
\begin{lstlisting}
# Dodaj novega VPN uporabnika
pivpn add

# Ogled aktivnih povezav
pivpn -c

# Ogled konfiguracije in QR koda
pivpn -qr

# Brisanje uporabnika
pivpn remove
\end{lstlisting}

\subsection{Konfiguracija in nastavitve}
Privzete pivpn nastavitve:
\begin{itemize}
    \item Podatkovna mapa: \verb|/etc/wireguard/|
    \item Datoteka strežnika: \verb|/etc/wireguard/wg0.conf|
    \item Konfiguracije VPN uporabnikov (QR kode, ključi): \verb|~/configs/| (domača mapa pivpn skripty)
    \item VPN subnet v našem sistemu: \texttt{10.4.103.0/24}
\end{itemize}

\subsection{Odpravljanje težav}
Če WireGuard ni dostopen ali se ne povezuje:
\begin{lstlisting}
# Preverka stanja vmesnika
ip link show wg0

# Preverka aktivnih povezav
wg show

# Resetiranje WireGuard vmesnika
sudo ip link delete wg0
sudo systemctl restart wg-quick@wg0
\end{lstlisting}

\subsection{wg-portal (WireGuard portal) in integracija z Active Directory}

\textbf{Kaj je wg-portal:} \texttt{wg-portal} je spletna aplikacija za upravljanje WireGuard konfiguracij, dovoljenj in omogočanje samopostrežnega ustvarjanja peer konfiguracij za uporabnike.

\textbf{Kje teče:} Aplikacija je dosegljiva na \texttt{http://192.168.11.106:8888}. Lokacije ostalih storitev so navedene v razdelku "DMZ strežniki" zgoraj.

\textbf{Kako se poveže z AD:}
\begin{itemize}
  \item Aplikacija uporablja LDAP bind račun za poizvedbe uporabnikov in preverjanje članstev v skupinah.
  \item Avtentikacija poteka z bindanjem (bind) in iskanjem vnosa, ki ustreza \texttt{sAMAccountName} ali \texttt{userPrincipalName}.
  \item Atribute AD lahko preslikamo, npr. uporabimo \texttt{userPrincipalName} kot \texttt{email}, kadar je \texttt{mail} prazen.
  \item Administratorske pravice se določijo preko članstva v AD skupini (npr. \texttt{Domain Admins}).
\end{itemize}

Pri namestitvi smo se v veliki meri držali vodnika: https://medium.com/@seeneeru/complete-guide-installing-and-configuring-wireguard-portal-on-linux-d23261027520 z nekaj lokalnimi modifikacijami prilagoditvami za naš okolje. Konfiguracija, ki smo jo uporabili v tem repozitoriju, je na voljo v datoteki \texttt{wgportal.yaml}.

\textbf{Opombe:}
\begin{itemize}
  \item Privzeta podatkovna shramba aplikacije je SQLite (\texttt{data/sqlite.db}). Če aplikacija poroča o napaki \texttt{unable to open database file}, preverite, da mapa \texttt{data} obstaja in ima pravilne pravice/lastništvo (npr. \texttt{/opt/wg-portal/data}).
  \item Primer ukazov za ustvarjanje bind uporabnika v AD (PowerShell):
\end{itemize}
\begin{lstlisting}
# Ustvari servisnega uporabnika (prilagodi geslo varno)
New-ADUser -Name "wgportal_bind" -SamAccountName wgportal_bind -AccountPassword
(ConvertTo-SecureString 'StrongP@ssw0rd' -AsPlainText -Force) -Enabled $true

# Po potrebi dodaj uporabnike v administratorsko skupino
Add-ADGroupMember -Identity "Domain Admins" -Members "zanadmin","pavlogal"
\end{lstlisting}

\textbf{Preizkusi in dnevniški pregledi:}
\begin{itemize}
  \item Testirajte LDAP poizvedbe z orodjem \texttt{ldapsearch} kjer je na voljo.
  \item Preverjajte dnevniške datoteke aplikacije za napake pri bindanju ali iskanju uporabnika.
  \item Če uporabniki ne vidijo konfiguracij, preverite, da so vklopljene nastavitve \newline  \verb|create_default_peer_on_login| in \verb|self_provisioning_allowed|.
\end{itemize}

\textbf{Opomba:} Po vsaki spremembi preveri sintakso in ponovno zaženi storitev, da se spremembe uveljavijo.


\section{REST}
Vse datoteke se nahajajo na sk11-dmz-linux-04 v home directoriju od zanzan04. Imamo dve rest storitvi in sicer Scriptum\_kp in library-api. V tem delu bom govoril o library-api za probleme s Scriptum\_kp pa se obrnite na github dokumentacijo projekta \url{https://github.com/zanzan731/Scriptum\_kp}.

\subsection{Arhitektura}
Komponente:
\begin{itemize}
  \item Node.js Express aplikacija (\texttt{server.js})
  \item SQLite baza (\texttt{library.db})
  \item TLS certifikati
  \item Docker in Docker Compose
  \item LDAP povezava proti Active Directory
\end{itemize}

\subsection{Konfiguracija}
Pomembne okoljske spremenljivke  (\texttt{server.js}):

\begin{lstlisting}
LDAP_URL=ldap://192.168.11.201:389
AD_DOMAIN=startup11.local
AD_BASE_DN=DC=startup11,DC=local
AD_LIBRARIAN_GROUP_DN=CN=LIBRARIAN,CN=Users,DC=startup11,DC=local
AD_ADMIN_USERS=zan_admin,other.user
\end{lstlisting}

Opis:
\begin{itemize}
  \item \texttt{LDAP\_URL} določa LDAP strežnik
  \item \texttt{AD\_BASE\_DN} določa osnovni DN
  \item \texttt{AD\_LIBRARIAN\_GROUP\_DN} določa skupino z dovoljenji
\end{itemize}

\subsection{Docker zagon}
Kljub temu da je možen lokalni zagon z \texttt{npm} (\texttt{npm install}, \texttt{npm run dev}), kar vam omogoča lokalno testiranje na \url{https://localhost:3443/} in \url{http://localhost:3000/} na serverju vedno buildamo to z dockerjem.

Build slike:
\begin{lstlisting}
docker build -t library-api-ad:latest .
\end{lstlisting}

Zagon kontejnerja:
\begin{lstlisting}
docker run -d \
  --name library-api-ad-test \
  -p 3000:3000 \
  -p 3443:3443 \
  -e LDAP_URL=ldap://192.168.11.201:389 \
  -e AD_DOMAIN=startup11.local \
  -e AD_BASE_DN=DC=startup11,DC=local \
  -e AD_LIBRARIAN_GROUP_DN=CN=LIBRARIAN,CN=Users,DC=startup11,DC=local \
  library-api-ad:latest
\end{lstlisting}

\subsection{Preverjanje delovanja}
Ko imate enkrat zagnano bi morala biti spletna stran dostopna preko vseh računalnikov znotraj omrežja na \url{https://192.168.11.104:3443/} in \url{http://192.168.11.104:3000/}.

\subsubsection{Javni GET zahtevek}
\begin{lstlisting}
curl -k -i https://192.168.11.104:3443/authors
curl --http2 -k -i https://192.168.11.104:3443/authors
curl -i http://192.168.11.104:3000/authors
\end{lstlisting}

\subsubsection{Vsebinsko pogajanje in podprti formati}
JSON:

\begin{lstlisting}
curl -k -H "Accept: application/json" \
https://192.168.11.104:3443/authors
\end{lstlisting}

XML:

\begin{lstlisting}
curl -k -H "Accept: application/xml" \
https://192.168.11.104:3443/authors
\end{lstlisting}

HTML:

\begin{lstlisting}
curl -k \
"https://192.168.11.104:3443/authors?format=html"
\end{lstlisting}

To je samo za authors imas pa tudi veliko drugih poti:
\begin{lstlisting}
GET /authors  - list all (public)
GET /authors/:id  - single author (public)
GET /authors/:id/books  - books by author (public)
POST /authors  - create (requires librarian or admin)
PUT /authors/:id  - update (requires librarian or admin)
DELETE /authors/:id  - delete (requires admin only)
GET /books  - list all, optional ?genre= filter (public)
GET /books/:id  - single book (public)
POST /books  - create (requires librarian or admin)
PUT /books/:id  - update (requires librarian or admin)
PATCH /books/:id/availability  - toggle availability (requires librarian or admin)
DELETE /books/:id  - delete (requires admin only)
\end{lstlisting}

\subsubsection{LDAP avtorizacija}
Rabis userja pod librarian ali admin za dodaten dostop.

\begin{itemize}
  \item branje (GET) avtorjev in knjig — javno dostopno brez prijave
  \item ustvarjanje (POST) avtorjev in knjig
  \item posodabljanje (PUT) avtorjev in knjig
  \item brisanje (DELETE) je rezervirano za admin uporabnike
\end{itemize}

Za windows PowerShell:
\begin{lstlisting}
[System.Net.ServicePointManager]::ServerCertificateValidationCallback = {$true}

$b64 = [Convert]::ToBase64String(
  [Text.Encoding]::ASCII.GetBytes(
    'username:password'
  )
)

$headers = @{
  Authorization = "Basic $b64"
  'Content-Type'='application/json'
}

$body = @{
  name = 'Test Author'
} | ConvertTo-Json

Invoke-RestMethod `
  -Uri 'https://192.168.11.104:3443/authors' `
  -Method Post `
  -Headers $headers `
  -Body $body
\end{lstlisting}

Za linux:
\begin{lstlisting}

curl -i -X POST http://192.168.11.104:3000/authors \
  -H "Content-Type: application/json" \
  -u username:password \
  -d '{"name":"Test Author from Linux"}'

curl -i -X DELETE http://192.168.11.104:3000/authors/5 \
  -u username:password
\end{lstlisting}

\subsubsection{Potrditev TLS certifikata}
Za certifikat:
\begin{lstlisting}
openssl s_client -connect 192.168.11.104:3443 -servername 192.168.11.104 -showcerts
curl -vkI https://192.168.11.104:3443/
\end{lstlisting}
Certifikati so self-signed zato če jim želiš zaupat jih rabiš naložiti v napravi:
\begin{lstlisting}
openssl s_client -connect 192.168.11.104:3443 -showcerts </dev/null \
  | sed -ne '/-BEGIN CERTIFICATE-/,/-END CERTIFICATE-/p' > /tmp/library-192-168-11-104.crt

sudo cp /tmp/library-192-168-11-104.crt /usr/local/share/ca-certificates/library-192-168-11-104.crt
sudo update-ca-certificates
\end{lstlisting}

\subsubsection{GraphQL}

Primer javnega GraphQL poizvedovanja:
\begin{lstlisting}
curl -k -X POST "https://192.168.11.104:32484/graphql" \
  -H "Content-Type: application/json" \
  -d '{"query":"{ authors { id name } }"}'
\end{lstlisting}

Primer zaščitene GraphQL mutacije z LDAP uporabnikom:
\begin{lstlisting}
$b64 = [Convert]::ToBase64String(
  [Text.Encoding]::ASCII.GetBytes('username:password')
)

$body = @{ query = 'mutation { createAuthor(name: "GraphQL Test") { id name } }' } | ConvertTo-Json

Invoke-RestMethod `
  -Uri 'https://192.168.11.104:32484/graphql' `
  -Method Post `
  -Headers @{ Authorization = "Basic $b64"; 'Content-Type'='application/json' } `
  -Body $body
\end{lstlisting}

\subsubsection{Scriptum\_kp}
Scriptum je dostopen na spletnem mestu \url{192.168.11.104:8080} in \url{88.200.24.241:8080} ne uporablja LDAP uporabnikov ampak uporabnike na MongoDB bazi. Frontend je Angular,, backend je REST API, Node.js, Express. MongoDB baza ni na tem serverju ampak je na zunanjih mongo DB strežnikih.

\section{Raft} 
Raft storitev poteka na treh serverjih 192.168.11.101, 192.168.11.102, 192.168.11.103. Vse spremembe so sihronizirane preko etcd kljucov, vsaka ima svojo lokalno bazo ki se sinhronizira. To je update Rest storitve saj ta sedaj pošilja svojo bazo RAFT podatkovni bazi.
Na serverju boste imeli veliko datotek:

\begin{itemize}
  \item \texttt{raft-api-integration.js} -  ta datoteka pomaga pri replikaciji in temu da se domenijo o vodji
  \item \texttt{sync-watcher.js} -  gleda etcd sync ključe in jih doda v localno library-ha.db za nek razlog etcd se ni znal sam zmenit za ključe
  \item \texttt{query-ha-db.js} -  node.js script da bere library-ha.db in sprinta autorje, uporabno za testiranje
  \item \texttt{test-graphql.js} -  test samo za lokalno da preveri če je možno POST request narediti na qraphql
  \item \texttt{create-author.js}, \texttt{manual-replicate.js} -  pomožne skripte za pomoč pri testiranju pisanja in repliciranja
\end{itemize}

\subsection{Branje baze}
Če želite manualno prebrati bazo predvsem za testiranje in če se vam zdi da niso sinhornizirane.
\begin{lstlisting}
cd ~/library-api
node query-ha-db.js
\end{lstlisting}
To vam bo izpisalo vsebino baze.




% Pretvorba med Markdown in LaTeX
%--------------------------------
% Ukazi, ki jih lahko uporabite za pretvorbo dokumentacije med formati.
% Zahteva: pandoc mora biti nameščen.
\section*{Pretvorba med Markdown in LaTeX}
\begin{itemize}
  \item LaTeX \textrightarrow{} Markdown (GitHub-flavored):
\begin{lstlisting}
pandoc -s vyos_documentation.tex -t gfm -o vyos_documentation.md --wrap=none
\end{lstlisting}

  \item Markdown \textrightarrow{} LaTeX:
\begin{lstlisting}
pandoc -s vyos_documentation.md -o vyos_documentation.tex
\end{lstlisting}
\end{itemize}

\end{document}

%